\documentclass{article}
\usepackage[utf8]{inputenc}
\usepackage{amsmath,amssymb}
\usepackage[most]{tcolorbox}
\usepackage{geometry}
\geometry{margin=2.5cm}

\newcommand{\step}[1]{\par\noindent\hangindent=3.5em\hangafter=1%
  \makebox[3.5em][l]{\textbf{#1.}}}
\newcommand{\steptag}[1]{\hfill{\small\textsf{[#1]}}}

\newtcolorbox{proofbox}[1]{
  colback=gray!5, colframe=gray!60,
  fonttitle=\bfseries, title={#1},
  breakable, enhanced,
  left=4pt, right=4pt, top=4pt, bottom=4pt,
  before skip=10pt, after skip=10pt
}

\begin{document}

\begin{proofbox}{Halo-robust height collapse: for an exact signed idempote...}
\step{1} Halo-robust height collapse: for an exact signed idempotent P with 0 < delta(P) <= 1/4, nonempty visible set W(P), and hidden top vertex v of height H, let sigma be the invisible mass of v, sigma\_g the halo-robust invisible mass (the positive coefficient mass v places on rows at ell-1 distance > tau/4 from conv W, tau = sqrt(delta)), and nu\_v the row negative mass; then H * (1 - sigma\_g) <= (sigma - sigma\_g) * tau/4 + nu\_v * (2 + 4*delta). \steptag{validated} \\[6pt]
\step{1.1} Set C=conv\{p\_w:w in W(P)\}, H\_j=dist\_1(p\_j,C), a\_j=P\_\{vj\}, a\_j\textasciicircum{}+=max(a\_j,0), a\_j\textasciicircum{}-=max(-a\_j,0), and tau=sqrt(delta). Let G=\{j:H\_j>tau/4\}, B=\{j:0<H\_j<=tau/4\}, and Z=\{j:H\_j=0\}. Then P\textasciicircum{}2=P gives p\_v=sum\_j a\_j p\_j; lem-mass-split gives sum\_j a\_j\textasciicircum{}+=1+nu\_v; sigma=sum\_\{j in G union B\} a\_j\textasciicircum{}+ and sigma\_g=sum\_\{j in G\} a\_j\textasciicircum{}+; hence sigma\_g<=sigma, sigma-sigma\_g=sum\_\{j in B\} a\_j\textasciicircum{}+, and sum\_\{j notin G\} a\_j\textasciicircum{}+ - sum\_j a\_j\textasciicircum{}-=1-sigma\_g. Also H\_j<=H for every row j, H\_j=0 on Z, H\_j<=tau/4 on B union Z, and every x in C satisfies ||x-p\_j||\_1<=2+4*delta for every row j. \steptag{validated} \\[6pt]
\step{1.1.1} For an exact signed idempotent P, P\textasciicircum{}2=P and a\_j=P\_\{vj\} imply row reproduction p\_v=sum\_j a\_j p\_j. The row sum condition gives sum\_j a\_j=1, and with a\_j=a\_j\textasciicircum{}+-a\_j\textasciicircum{}- and nu\_v=sum\_j a\_j\textasciicircum{}- for row v, lem-mass-split gives sum\_j a\_j\textasciicircum{}+=1+nu\_v. \steptag{validated} \\[4pt]
\step{1.1.2} By def-invisible-mass and the halo-robust sigma\_g definition in the contract, sigma is the sum of a\_j\textasciicircum{}+ over rows with H\_j>0 and sigma\_g is the sum of a\_j\textasciicircum{}+ over rows with H\_j>tau/4. With G=\{H\_j>tau/4\}, B=\{0<H\_j<=tau/4\}, and Z=\{H\_j=0\}, this gives sigma=sum\_\{G union B\}a\_j\textasciicircum{}+, sigma\_g=sum\_G a\_j\textasciicircum{}+, sigma\_g<=sigma, sigma-sigma\_g=sum\_B a\_j\textasciicircum{}+, and sum\_\{j notin G\}a\_j\textasciicircum{}+-sum\_j a\_j\textasciicircum{}-=1-sigma\_g by node 1.1.1. \steptag{validated} \\[4pt]
\step{1.1.3} By def-height, v is a hidden top vertex with H=dist\_1(p\_v,C) equal to the maximum row distance from C, so H\_j<=H for every row j; by the definitions of B and Z, H\_j<=tau/4 on B union Z and H\_j=0 on Z. By the row-geometry clause of def-signed-idempotent, pairwise row distances are at most 2+4*delta; since every x in C is a convex combination of visible rows, convexity of the l1 norm gives ||x-p\_j||\_1<=2+4*delta for every x in C and every row j. \steptag{validated} \\[4pt]
\step{1.2} If sigma\_g>=1, the desired inequality follows immediately: H>=0 gives H*(1-sigma\_g)<=0, while sigma-sigma\_g>=0, tau>=0, nu\_v>=0, and 2+4*delta>=0 make the right side nonnegative. \steptag{validated} \\[4pt]
\step{1.3} Assume sigma\_g<1 and define q=(p\_v-sum\_\{j in G\} a\_j\textasciicircum{}+ p\_j)/(1-sigma\_g). Since p\_v=sum\_\{j in G\} a\_j\textasciicircum{}+p\_j+(1-sigma\_g)q, the coefficients a\_j\textasciicircum{}+ on G are nonnegative and have sum sigma\_g<1, and H\_j<=H=dist\_1(p\_v,C) for j in G, lem-residual-lower gives H<=dist\_1(q,C); multiplying by 1-sigma\_g gives H*(1-sigma\_g)<= (1-sigma\_g)*dist\_1(q,C). \steptag{validated} \\[4pt]
\step{1.4} Under sigma\_g<1, the same q has numerator sum\_\{j notin G\} a\_j\textasciicircum{}+ p\_j - sum\_j a\_j\textasciicircum{}- p\_j and denominator m=1-sigma\_g>0. Applying lem-residual-upper with positive coefficients a\_j\textasciicircum{}+ for j notin G, negative coefficients a\_j\textasciicircum{}- for all j, and D\_j=2+4*delta gives (1-sigma\_g)*dist\_1(q,C) <= sum\_\{j notin G\} a\_j\textasciicircum{}+ H\_j + nu\_v*(2+4*delta) <= (sigma-sigma\_g)*tau/4 + nu\_v*(2+4*delta). \steptag{validated} \\[6pt]
\step{1.4.1} When sigma\_g<1 and q=(p\_v-sum\_\{j in G\}a\_j\textasciicircum{}+p\_j)/(1-sigma\_g), row reproduction gives q=(sum\_\{j notin G\}a\_j\textasciicircum{}+p\_j - sum\_j a\_j\textasciicircum{}-p\_j)/(1-sigma\_g). Node 1.1.2 gives m=sum\_\{j notin G\}a\_j\textasciicircum{}+-sum\_j a\_j\textasciicircum{}-=1-sigma\_g>0. \steptag{validated} \\[4pt]
\step{1.4.2} The hypotheses of lem-residual-upper hold for C, b\_j=a\_j\textasciicircum{}+ on j notin G, c\_j=a\_j\textasciicircum{}- on all rows, p\_j the corresponding rows, r\_j the corresponding negative-source rows, m=1-sigma\_g, and D\_j=2+4*delta: all coefficients are nonnegative, q has the required quotient form, and node 1.1.3 supplies ||x-r\_j||\_1<=D\_j for every x in C. \steptag{validated} \\[4pt]
\step{1.4.3} By lem-residual-upper, (1-sigma\_g)*dist\_1(q,C) <= sum\_\{j notin G\}a\_j\textasciicircum{}+H\_j + nu\_v*(2+4*delta). Since notin G is B union Z, node 1.1.3 gives H\_j<=tau/4 on B and H\_j=0 on Z, while node 1.1.2 gives sum\_B a\_j\textasciicircum{}+=sigma-sigma\_g; therefore sum\_\{j notin G\}a\_j\textasciicircum{}+H\_j <= (sigma-sigma\_g)*tau/4. \steptag{validated} \\[6pt]
\step{1.4.3.1} Under the sigma\_g<1 branch and the q defined in node 1.3, the quotient form needed for lem-residual-upper is available inside this node: node 1.1.1 gives p\_v=sum\_j(a\_j\textasciicircum{}+-a\_j\textasciicircum{}-)p\_j, so subtracting sum\_\{j in G\}a\_j\textasciicircum{}+p\_j from the definition (1-sigma\_g)q=p\_v-sum\_\{j in G\}a\_j\textasciicircum{}+p\_j gives (1-sigma\_g)q=sum\_\{j notin G\}a\_j\textasciicircum{}+p\_j-sum\_j a\_j\textasciicircum{}-p\_j. Node 1.1.2 gives sum\_\{j notin G\}a\_j\textasciicircum{}+-sum\_j a\_j\textasciicircum{}-=1-sigma\_g, and the current branch has sigma\_g<1, hence m=1-sigma\_g>0. \steptag{validated} \\[4pt]
\step{1.4.3.2} For this application of lem-residual-upper, take C=conv W, the positive coefficients b\_j=a\_j\textasciicircum{}+ over indices j notin G, the negative coefficients c\_k=a\_k\textasciicircum{}- over all row indices k, the corresponding points p\_j and r\_k to be those rows of P, and D\_k=2+4*delta. All coefficients are nonnegative by definition of positive and negative parts. The quotient hypothesis holds because row reproduction gives p\_v=sum\_j(a\_j\textasciicircum{}+-a\_j\textasciicircum{}-)p\_j and q is defined by (1-sigma\_g)q=p\_v-sum\_\{j in G\}a\_j\textasciicircum{}+p\_j, so q=(sum\_\{j notin G\}a\_j\textasciicircum{}+p\_j-sum\_k a\_k\textasciicircum{}-p\_k)/(1-sigma\_g); node 1.1.2 gives m=sum\_\{j notin G\}a\_j\textasciicircum{}+-sum\_k a\_k\textasciicircum{}-=1-sigma\_g, and the branch sigma\_g<1 gives m>0. Finally D\_k>=0 and node 1.1.3 gives ||x-r\_k||\_1<=D\_k for every x in C. Thus lem-residual-upper yields (1-sigma\_g)*dist\_1(q,C)<=sum\_\{j notin G\}a\_j\textasciicircum{}+ dist\_1(p\_j,C)+sum\_k a\_k\textasciicircum{}-(2+4*delta)=sum\_\{j notin G\}a\_j\textasciicircum{}+H\_j+nu\_v*(2+4*delta). \steptag{validated} \\[4pt]
\step{1.4.3.3} The remaining sum is bounded by local halo bookkeeping: node 1.4.3.3.1 gives the disjoint partition \{j:notin G\}=B union Z, the bounds H\_j<=tau/4 on B and H\_j=0 on Z, and the mass identity sum\_\{j in B\}a\_j\textasciicircum{}+=sigma-sigma\_g. Therefore, as shown in node 1.4.3.3.2, nonnegativity of a\_j\textasciicircum{}+ and H\_j gives sum\_\{j notin G\}a\_j\textasciicircum{}+H\_j=sum\_\{j in B\}a\_j\textasciicircum{}+H\_j+sum\_\{j in Z\}a\_j\textasciicircum{}+H\_j <= (tau/4)*sum\_\{j in B\}a\_j\textasciicircum{}+=(sigma-sigma\_g)*tau/4. \steptag{validated} \\[6pt]
\step{1.4.3.3.1} Local halo partition and mass identity, with the setup fixed explicitly: let C=C\_W=conv\{p\_w:w in W(P)\}, H\_j=dist\_1(p\_j,C\_W), a\_j=P\_\{vj\}, a\_j\textasciicircum{}+=max(a\_j,0), tau=sqrt(delta), G=\{j:H\_j>tau/4\}, B=\{j:0<H\_j<=tau/4\}, and Z=\{j:H\_j=0\}. Since H\_j is a distance to C\_W, H\_j>=0 for every j, so \{j:notin G\}=B union Z disjointly. By def-invisible-mass, sigma is the sum of a\_j\textasciicircum{}+ over rows with dist\_1(p\_j,C\_W)>0, equivalently H\_j>0. By the contract definition of the halo-robust invisible mass, sigma\_g is the sum of a\_j\textasciicircum{}+ over rows with dist\_1(p\_j,C\_W)>tau/4, equivalently H\_j>tau/4, i.e. over G. Since \{H\_j>0\}=G union B disjointly, sigma=sum\_\{j in G\}a\_j\textasciicircum{}+ + sum\_\{j in B\}a\_j\textasciicircum{}+ and sigma\_g=sum\_\{j in G\}a\_j\textasciicircum{}+; hence sum\_\{j in B\}a\_j\textasciicircum{}+=sigma-sigma\_g. The definitions of B and Z also give H\_j<=tau/4 on B and H\_j=0 on Z. \steptag{validated} \\[6pt]
\step{1.4.3.3.1.1} Explicit setup bridge for the verifier objection: in node 1.4.3.3.1 the symbol C is not arbitrary; it is C\_W=conv\{p\_w:w in W(P)\} from def-visible-set and def-height. The coefficients are the row-v coefficients a\_j=P\_\{vj\}, with positive parts a\_j\textasciicircum{}+=max(a\_j,0). Therefore def-invisible-mass applies exactly as sigma=sum\_\{j:dist\_1(p\_j,C\_W)>0\}a\_j\textasciicircum{}+, while the contract defines sigma\_g=sum\_\{j:dist\_1(p\_j,C\_W)>tau/4\}a\_j\textasciicircum{}+. Since H\_j=dist\_1(p\_j,C\_W), the sets \{H\_j>0\} and \{H\_j>tau/4\} are respectively G union B and G; hence sum\_\{j in B\}a\_j\textasciicircum{}+=sigma-sigma\_g. Nonnegativity of distance gives the disjoint partition \{j:notin G\}=B union Z, and the definitions of B and Z give H\_j<=tau/4 on B and H\_j=0 on Z. \steptag{validated} \\[4pt]
\step{1.4.3.3.2} With H\_j=dist\_1(p\_j,C), G=\{j:H\_j>tau/4\}, B=\{j:0<H\_j<=tau/4\}, and Z=\{j:H\_j=0\}, nonnegativity of distances gives the disjoint partition \{j:notin G\}=B union Z. By def-invisible-mass, sigma is the sum of a\_j\textasciicircum{}+ over rows with H\_j>0, and by the contract sigma\_g is the sum of a\_j\textasciicircum{}+ over rows with H\_j>tau/4, so sigma-sigma\_g=sum\_\{j in B\}a\_j\textasciicircum{}+. Also H\_j=0 on Z and H\_j<=tau/4 on B. Since a\_j\textasciicircum{}+>=0, sum\_\{j notin G\}a\_j\textasciicircum{}+H\_j=sum\_\{j in B\}a\_j\textasciicircum{}+H\_j+sum\_\{j in Z\}a\_j\textasciicircum{}+H\_j <= (tau/4)*sum\_\{j in B\}a\_j\textasciicircum{}+=(sigma-sigma\_g)*tau/4. \steptag{validated} \\[4pt]
\step{1.4.4} Boundary-inclusive halo bookkeeping for the second inequality: set C=C\_W=conv\{p\_w:w in W(P)\}, H\_j=dist\_1(p\_j,C\_W), a\_j=P\_\{vj\}, a\_j\textasciicircum{}+=max(a\_j,0), tau=sqrt(delta), G=\{j:H\_j>tau/4\}, N=\{j:0<H\_j<=tau/4\}, and Z=\{j:H\_j=0\}. Since H\_j is a distance, H\_j>=0 and hence \{j:notin G\}=N union Z disjointly, including the boundary H\_j=tau/4 in N. By def-invisible-mass, sigma=sum\_\{H\_j>0\} a\_j\textasciicircum{}+; by the contract definition of halo-robust invisible mass, sigma\_g=sum\_\{H\_j>tau/4\} a\_j\textasciicircum{}+=sum\_\{j in G\}a\_j\textasciicircum{}+. Therefore sigma-sigma\_g=sum\_\{j in N\}a\_j\textasciicircum{}+. On Z the summand a\_j\textasciicircum{}+ H\_j is zero, and on N one has H\_j<=tau/4, so nonnegativity of a\_j\textasciicircum{}+ gives sum\_\{j notin G\} a\_j\textasciicircum{}+ H\_j=sum\_\{j in N\}a\_j\textasciicircum{}+H\_j+sum\_\{j in Z\}a\_j\textasciicircum{}+H\_j <= (tau/4) sum\_\{j in N\}a\_j\textasciicircum{}+=(sigma-sigma\_g)*tau/4. Thus the boundary rows are accounted for and the displayed final inequality in node 1.4 follows when this is combined with the lem-residual-upper bound already established for the same q and G. \steptag{validated} \\[4pt]
\step{1.5} Combining the two cases proves the contract: the sigma\_g>=1 case is node 1.2, and in the sigma\_g<1 case nodes 1.3 and 1.4 chain to H*(1-sigma\_g) <= (sigma-sigma\_g)*tau/4 + nu\_v*(2+4*delta). \steptag{validated} \\[6pt]
\step{1.5.1} Conditional combination step, available only after nodes 1.2, 1.3, and 1.4 are validated: if sigma\_g>=1, node 1.2 gives exactly the desired inequality. If sigma\_g<1, node 1.3 gives H*(1-sigma\_g) <= (1-sigma\_g)*dist\_1(q,C), and node 1.4 gives (1-sigma\_g)*dist\_1(q,C) <= (sigma-sigma\_g)*tau/4 + nu\_v*(2+4*delta); transitivity gives H*(1-sigma\_g) <= (sigma-sigma\_g)*tau/4 + nu\_v*(2+4*delta). Since exactly one of sigma\_g>=1 and sigma\_g<1 holds, the two validated cases prove the contract. \steptag{validated} \\[4pt]
\end{proofbox}

\end{document}
